\documentclass[../main.tex]{subfiles}
\begin{document}

Энэ бүлэгт системийг хэрхэн туршсан арга зүй болон хөгжүүлэлтийн явцад
тулгарсан гурван томоохон асуудлыг асуудал $\rightarrow$ оношилгоо
$\rightarrow$ шийдэл $\rightarrow$ баталгаажуулалт гэсэн дарааллаар
тайлбарлав. Эдгээр асуудал бүр анх төлөвлөөгүй боловч бодит хэрэглээнд
гарсан бөгөөд шийдвэрлэх явц нь дадлагын хамгийн их сургамжтай хэсэг байв.

\section{Туршилтын арга зүй}

Системийг гурван түвшинд туршив. Нэгдүгээрт, төслийн хэмжээнд нэгж тест
бичихийн оронд бүрэн integration smoke тест (\texttt{scripts/smoke-test.sh})
бичсэн бөгөөд энэ нь ажиллаж буй стекийн эсрэг дараах бүрэн урсгалыг
автоматаар шалгана: шинэ хэрэглэгч бүртгэх $\rightarrow$ нэвтрээгүй хүсэлт
401 авахыг батлах $\rightarrow$ текст буулгах $\rightarrow$ баримт
\texttt{ready} болтол хүлээх $\rightarrow$ чат асуулт илгээх $\rightarrow$
хариулт эх сурвалжтай ирснийг шалгах $\rightarrow$ харилцан ярианы түүх
хадгалагдсаныг батлах. Хоёрдугаарт, өөрчлөлт бүрийн дараа бүх workspace-д
\texttt{pnpm typecheck} болон \texttt{pnpm build} ажиллуулж хэлний түвшний
алдааг барьдаг. Гуравдугаарт, хайлтын параметр зэрэг UI-ээс хамааралгүй
шинж чанаруудыг \texttt{curl}-ээр шууд, харин мобайл аппыг бодит утсан дээр
хоёр удаагийн туршилтын үе шаттайгаар гараар туршив.

\section{Асуудал 1: Хэл дамнасан хайлтын чанар}
\label{sec:problem-crosslingual}

\textbf{Асуудал.} Кирилл монголоор бичигдсэн баримтаас латин үсгээр (жишээ
нь англиар) асуулт асуухад систем баримтад огт байхгүй, өөр хэл дээрх
хариулт өгч байв --- өөрөөр хэлбэл RAG системийн үндсэн зорилго болох
«баримтад суурилсан хариулт» зөрчигдөж байв.

\textbf{Оношилгоо.} Шалтгаан нь хоёр давхаргад зэрэг байсан бөгөөд нэг нь
нөгөөгөө далдалж байв. Хайлтын давхаргад: \texttt{gemini-embedding-001}
загвар латин асуулт болон кирилл хэсгүүдийг embedding огторгуйн өөр өөр
мужид байрлуулдаг тул хоорондын косинус төстэй байдал ойролцоогоор
\textbf{0.35} гарч, анхны \texttt{threshold=0.7} босгыг давахгүй бүх
кандидат шүүгдэн хаягдаж, context хоосон үлдэж байв. Үүсгэлтийн давхаргад:
хуучин системийн заавар зөвхөн «баримтаас хариул» гэж заасан тул context
хоосон ирэхэд загвар өөрийн сургалтын мэдлэгээс, асуултын хэлээр (англиар)
хариулж байв. Аль нэгийг нь дангаар нь засахад асуудал бүрэн арилахгүй
байсан нь оношилгоог төвөгтэй болгосон.

\textbf{Шийдэл.} Гурван өөрчлөлт хийв: (1) \texttt{threshold}-ыг $0.7$-оос
$0.1$ болгож бууруулав --- энэ нь хэл дамнасан хайлтыг нээж өгсөн түр
шийдэл бөгөөд урт хугацаанд асуултыг баримтын хэл рүү орчуулж хайдаг болгох
нь зөв шийдэл гэж баримтжуулсан; (2) MMR-ийн $\lambda$-г $0.5$ болгож олон
талт байдлын жинг нэмэгдүүлэв; (3) системийн зааврыг бүрэн шинэчилж,
хариултын хэлийг асуултын хэлээс бус \emph{баримтын} хэлнээс хамааруулах
дүрэм оруулав (5.2 хэсэгт иш татсан).

\textbf{Баталгаажуулалт.} Threshold-ын нөлөөг тусгаарлаж шалгахын тулд
\texttt{/chat}-ын параметрээр хоёр туйлын утга өгч туршив:
\texttt{threshold=0.95} үед систем контекст олохгүй тул «мэдээлэл
олдсонгүй» гэсэн шударга хариу (0 эх сурвалжтай) буцаасан бол
\texttt{threshold=0.1} үед латин асуултад кирилл баримтаас 5 эх сурвалжтай,
баримтад үндэслэсэн кирилл хариулт өгөв. Энэ туршилт мөн хайлтын
тохиргооны UI (5.4 хэсэг) зөв дамжиж буйг давхар баталсан.

\section{Асуудал 2: Том хэмжээний текст оруулах}
\label{sec:problem-large-paste}

\textbf{Асуудал.} 50{,}000 тэмдэгттэй кирилл текст буулгахад баримт хэзээ ч
\texttt{ready} болохгүй, ямар ч шалтгаан заалгүйгээр \texttt{failed} төлөвт
ордог байв. Түүнчлэн баримтжуулсан дээд хэмжээ болох 500{,}000 тэмдэгттэй
текст серверт огт хүрэлгүй \texttt{413 Payload Too Large} алдаа авч байв.

\textbf{Оношилгоо.} API логоос эхний \texttt{batchEmbedContents} хүсэлт
шууд \texttt{429 RESOURCE\_EXHAUSTED} буцааж байгааг илрүүлэв. Хуучин
embedder нэг хүсэлтэд 100 хүртэл хэсэг ($\approx$23 мянган токен) багцалж
илгээдэг байсан бол Gemini-ийн үнэгүй түвшин минутад ойролцоогоор 20--25
мянган токены квот хэрэгжүүлдэг нь туршилтаар тогтоогдов: 6 мянган токентой
гурван багц дараалан амжилттай өнгөрч, дөрөв дэх нь 429 авдаг бол 23 мянган
токентой дан багц шууд 429 авдаг. Давтан оролдлого байгаагүй тул эхний 429
бүхэл pipeline-ыг унагаж байв. Харин 413 алдааны шалтгаан нь кодировкод
байв: кирилл тэмдэгт UTF-8-д 2 байт эзэлдэг тул 500 мянган тэмдэгт
$\approx$1~MB болж Fastify-ийн анхдагч 1~MB \texttt{bodyLimit}-ээс хальж
байв --- өөрөөр хэлбэл баримтжуулсан дээд хэмжээ монгол текстийн хувьд
огт хүрэх боломжгүй байжээ.

\textbf{Шийдэл.} Embedder-ийг бүрэн дахин бичив (5.1 хэсэгт тайлбарласан):
багц бүрийг 15 хэсэг буюу 6 мянган токеноор хязгаарлаж, багц хооронд 3
секунд түр зогсож, 429/5xx алдаанд exponential backoff-оор 8 хүртэл удаа
дахин оролддог болгов. \texttt{bodyLimit}-ыг 4~MB болгож өсгөв. Мөн чухал
нэг сайжруулалт: алдааны шалтгааныг \texttt{documents.error\_message}
баганад (migration 003) хадгалж, вэб болон мобайл интерфэйс дээр
хэрэглэгчид харуулдаг болгосноор «чимээгүй уналт» дахин гарахааргүй болов.

\textbf{Баталгаажуулалт.} 50 мянган тэмдэгттэй текст 50 хэсэг / 5 багц болж,
дараалсан таван 429 алдааг backoff-оор даван ойролцоогоор 95 секундэд
\texttt{ready} болов. 500 мянган тэмдэгттэй текст 413 авахгүй хүлээн
авагдаж, 492 хэсэг / 41 багц болон тогтвортой хэмнэлтэйгээр боловсрогдож
байгаа нь ажиглагдсан (бүрэн гүйцэтгэл нь туршилтын орчны Docker-ийн
саатлаар тасалдсан тул механизм нь 50 мянгын туршилтаар батлагдсанд
тооцов). Энэ оношилгооны явцад бас нэг сургамж гарсан: Docker Desktop-ийн
bind mount засварласан кодын оронд хуучин кодыг үйлчилсээр байх тохиолдол
илэрсэн тул «боломжгүй» алдааг мөшгихийн өмнө контейнер доторх кодыг
\texttt{docker exec} + \texttt{grep}-ээр шалгах дүрэм тогтоов.

\section{Асуудал 3: Мобайл төхөөрөмжийн boot loop}
\label{sec:problem-bootloop}

\textbf{Асуудал.} Бодит утсан дээрх анхны туршилтаар апп огт ашиглагдахгүй
байв: танилцуулга дуусаад «эхлэх» товч дарахад 3 орчим минут ачаалж, дараа
нь дахин танилцуулга руу буцдаг хязгааргүй давталт үүсэв. Эмулятор дээр
энэ асуудал огт илэрдэггүй байсан нь оношилгоог хүндрүүлсэн.

\textbf{Оношилгоо.} Гурван алдаа давхарласан байв: (1)
\texttt{EXPO\_PUBLIC\_API\_URL} нь Android эмуляторын тусгай хаяг болох
\texttt{10.0.2.2}-т заасан байсан нь бодит утаснаас хүрэх боломжгүй; (2)
session шалгах \texttt{getSession()} дуудлагад timeout байгаагүй тул
хүрэхгүй хаяг руу минутаар дүүжигнэж байв; (3) алдааны \texttt{catch}
блок ямар ч алдааг «онбординг хийгээгүй» гэж тайлбарлан \texttt{/onboarding}
руу чиглүүлдэг байсан нь давталтыг үүсгэв.

\textbf{Шийдэл.} Гурвууланг нь зэрэг засав: API хаягийг олон кандидатаас
зэрэгцээ liveness шалгалтаар автоматаар олдог \texttt{api-base.ts} модуль
бичиж (5.5 хэсэгт тайлбарласан), session шалгалтыг 5 секундын timeout-той
\texttt{Promise.race} дотор оруулж, онбординг дууссан хэрэглэгчийн алдааг
\texttt{/login} руу чиглүүлдэг болгов --- ямар ч тохиолдолд онбординг руу
буцах зам байхгүй болсон.

\textbf{Баталгаажуулалт.} Дараагийн буюу хоёр дахь утасны туршилтаар boot
loop бүрэн арилж, бүртгэл $\rightarrow$ нэвтрэлт $\rightarrow$ баримт
байршуулалт $\rightarrow$ чат гэсэн бүрэн урсгал бодит төхөөрөмж дээр
ажилласан. Энэ туршилтаас гарсан таван жижиг UX асуудал (жишээ нь чатын
урт түүхэн дээр touch scroll ажиллахгүй болох, онбордингийн «дахиж
харуулахгүй» сонголт үйлчлэхгүй байх) мөн адил оношлогдож засагдсан.

\section{Гүйцэтгэлийн хэмжилтүүд}

Баримт оруулах хоолойн гүйцэтгэлийг янз бүрийн хэмжээтэй кирилл текстээр
хэмжсэн дүнг Хүснэгт~\ref{tab:perf}-д нэгтгэв. Хугацааны дийлэнх нь Gemini
API-ийн үнэгүй түвшний квотод зориулсан хүлээлт тул төлбөрт түвшинд
шилжвэл эдгээр тоо эрс багасна гэдгийг тэмдэглэх нь зүйтэй.

\renewcommand{\arraystretch}{1.3}
\begin{longtable}{|>{\raggedright\arraybackslash}m{3.4cm}
                  |>{\centering\arraybackslash}m{2cm}
                  |>{\centering\arraybackslash}m{2cm}
                  |>{\centering\arraybackslash}m{3cm}
                  |>{\raggedright\arraybackslash}m{4.2cm}|}
\caption{Баримт оруулах хоолойн гүйцэтгэл (кирилл текст)}
\label{tab:perf}
\\ \hline
\textbf{Оролт} & \textbf{Chunk} & \textbf{Багц} & \textbf{Хугацаа} &
\textbf{Тэмдэглэл} \\ \hline
\endfirsthead
\hline
\textbf{Оролт} & \textbf{Chunk} & \textbf{Багц} & \textbf{Хугацаа} &
\textbf{Тэмдэглэл} \\ \hline
\endhead

Богино текст ($\sim$3 мянган тэмдэгт) & 2--3 & 1 & $<$10 сек &
Нэг багцад багтана, хүлээлтгүй \\ \hline

50 мянган тэмдэгт & 50 & 5 & $\sim$95 сек &
Таван дараалсан 429-ийг backoff-оор даван амжилттай \texttt{ready} \\ \hline

500 мянган тэмдэгт & 492 & 41 & $\sim$25--30 мин (тооцоолол) &
413-гүй хүлээн авагдаж, тогтвортой хэмнэлээр боловсрогдсон; механизм 50
мянгын туршилтаар батлагдсан \\ \hline
\end{longtable}

Чатын хариултын хугацаа сонгосон загвараас хамаарч 3--10 секундын хооронд
хэлбэлзэж байсан бөгөөд үүний багагүй хэсэг нь Gemini-ийн generation
дуудлага өөрөө юм. Smoke тест бүрэн урсгалаа тогтмол давдаг, бүх workspace-ийн
typecheck болон build алдаагүй байгаа нь системийн одоогийн тогтвортой
байдлын үндсэн үзүүлэлт болж байна.

\end{document}
